% Fig. 2a — charge-transport mechanism schematic.
% Quantitative panels b–d are owned by the data pipeline; this fixture draws
% only the reader-facing mechanism and the two-window readout they quantify.
\documentclass[border=0pt]{standalone}
\usepackage{polymer-paper-preamble}
\usetikzlibrary{arrows.meta,calc,decorations.pathmorphing}

\begin{document}
\resizebox{163.8mm}{!}{%
\begin{tikzpicture}[
  every node/.style={font=\sffamily\fontsize{6.0}{7.1}\selectfont},
  panelLetter/.style={font=\sffamily\bfseries\fontsize{8}{9.6}\selectfont,
    text=cGray!92!black},
  panelTitle/.style={font=\sffamily\bfseries\fontsize{6.6}{7.8}\selectfont,
    text=cGray!82!black},
  columnTitle/.style={font=\sffamily\bfseries\fontsize{6.8}{8.1}\selectfont,
    text=cGray!86!black},
  labelStd/.style={font=\sffamily\fontsize{5.7}{6.8}\selectfont,
    text=cGray!84!black},
  labelMute/.style={font=\sffamily\itshape\fontsize{5.25}{6.3}\selectfont,
    text=cGray!63!black},
  wire/.style={cGray!72!black,line width=0.65pt,line cap=round,line join=round},
  electrode/.style={fill=cGray!20,draw=cGray!74!black,line width=0.65pt},
  film/.style={fill=cAmber!10,draw=cAmber!72!black,line width=0.75pt},
  conventionalFilm/.style={fill=cBlue!6,draw=cBlue!58!black,line width=0.75pt},
  trapBlue/.style={circle,fill=cBlue!72,draw=cBlue!60!black,line width=0.3pt,
    inner sep=0pt,minimum size=1.7mm},
  trapRed/.style={circle,fill=cRed!78,draw=cRed!60!black,line width=0.3pt,
    inner sep=0pt,minimum size=1.7mm},
  trapAmber/.style={circle,fill=cAmber!88,draw=cRed!58!black,line width=0.3pt,
    inner sep=0pt,minimum size=1.75mm},
  drainArrow/.style={-{Stealth[length=3.4pt,width=2.5pt]},cBlue!72!black,
    line width=0.62pt},
  storyArrow/.style={-{Stealth[length=3.4pt,width=2.4pt]},cGray!62!black,
    line width=0.52pt},
  comparisonArrow/.style={{Stealth[length=3.1pt,width=2.2pt]}-{Stealth[length=3.1pt,width=2.2pt]},
    cGray!58!black,line width=0.48pt},
  axisArrow/.style={-{Stealth[length=3.2pt,width=2.2pt]},cGray!68!black,
    line width=0.55pt},
  traceEarly/.style={cRed!76!black,line width=0.9pt,line cap=round},
  traceLate/.style={cRed!70!black,line width=0.9pt,line cap=round,
    dash pattern=on 2pt off 1.3pt},
]

% Preserve the allocated Fig. 2a row height while allowing the open canvas to
% carry less decorative geometry than the former full-height separators.
\path[use as bounding box] (0.00,0.38) rectangle (16.38,5.02);

% A single, open mechanism strip. Measurement feeds a parallel material
% comparison; both material states feed the common qualitative readout.
\node[panelLetter,anchor=west] at (0.14,4.88) {a};

% Reading-order cues stay on a shared mechanism rail; the paired arrow marks
% a comparison, not a sequential conversion from one dielectric into another.
\draw[storyArrow] (3.38,3.62)--(4.24,3.62);
\draw[comparisonArrow] (7.30,3.62)--(8.18,3.62);
\node[labelMute,anchor=south] at (7.74,3.68) {material contrast};
\draw[storyArrow] (11.46,3.62)--(12.42,3.62);

% --------------------------------------------------------------------------
% Measurement
% --------------------------------------------------------------------------
\node[columnTitle,anchor=north] at (2.02,4.72) {Measurement};
\node[labelMute,anchor=north] at (2.02,4.30) {MIM stack};

% A compact vertical MIM stack: the film is the amber middle layer, not a
% second electrode or a free-standing sample.
\fill[electrode] (0.78,2.98) rectangle (3.22,3.28);
\fill[film] (0.78,2.10) rectangle (3.22,2.98);
\fill[electrode] (0.78,1.80) rectangle (3.22,2.10);
\node[labelStd,text=cRed!70!black,anchor=west] at (0.18,1.46)
  {$E=15\,\mathrm{MV\,m^{-1}}$};
\node[labelMute,anchor=north] at (2.02,0.80) {25\,\%RH chamber};

% --------------------------------------------------------------------------
% Conventional dielectric
% --------------------------------------------------------------------------
\node[columnTitle,anchor=north] at (5.72,4.72) {Conventional dielectric};
\node[labelMute,anchor=north] at (5.72,4.30) {finite shallow sites};

\fill[conventionalFilm] (4.40,2.12) rectangle (7.04,3.40);
\draw[wire] (4.40,2.12)--(7.04,2.12);
\draw[wire] (4.40,3.40)--(7.04,3.40);
\node[trapBlue] at (5.05,2.78) {};
\node[trapBlue] at (5.72,2.43) {};
\node[trapBlue] at (6.40,2.95) {};
\draw[drainArrow] (5.05,2.69)--(5.05,1.68);
\draw[drainArrow] (5.72,2.34)--(5.72,1.55);
\draw[drainArrow] (6.40,3.04)--(6.40,1.68);
\node[labelStd,text=cBlue!66!black,anchor=north] at (5.72,1.18)
  {charge drains};

% --------------------------------------------------------------------------
% Sulfur-rich copolymer
% --------------------------------------------------------------------------
\node[columnTitle,anchor=north] at (9.78,4.72) {Sulfur-rich copolymer};
\node[labelMute,anchor=north] at (9.78,4.30) {distributed trap population};

\fill[film] (8.38,2.12) rectangle (11.18,3.40);
\draw[wire] (8.38,2.12)--(11.18,2.12);
\draw[wire] (8.38,3.40)--(11.18,3.40);
% Non-periodic host traces make the breadth visible without recreating Fig1's
% energy diagram or implying a calibrated density distribution.
\draw[cGray!48!black,line width=0.48pt,line cap=round]
  (8.58,2.88) .. controls (8.88,2.56) and (9.12,3.10) .. (9.48,2.62);
\draw[cAmber!62!black,line width=0.55pt,line cap=round]
  (9.46,2.34) .. controls (9.82,2.14) and (10.10,2.84) .. (10.44,2.58);
\draw[cGray!48!black,line width=0.48pt,line cap=round]
  (10.30,3.08) .. controls (10.62,2.76) and (10.80,3.20) .. (11.02,2.68);
\node[trapAmber] at (8.88,2.76) {};
\node[trapRed] at (9.36,2.56) {};
\node[trapAmber] at (9.88,2.44) {};
\node[trapRed] at (10.20,2.86) {};
\node[trapAmber] at (10.60,2.98) {};
\node[trapRed] at (10.96,2.72) {};
% The occupied sites remain inside the film and have no outward drain arrows;
% that restrained contrast carries retention without a decorative halo.
\node[labelStd,text=cRed!68!black,anchor=north] at (9.78,1.18)
  {charge retained};

% --------------------------------------------------------------------------
% What we read
% --------------------------------------------------------------------------
\node[columnTitle,anchor=north] at (14.34,4.72) {What we read};
\node[labelMute,anchor=north] at (14.34,4.30) {one decay, two windows};

\draw[axisArrow] (12.92,1.95)--(12.92,3.64);
\draw[axisArrow] (12.92,1.95)--(15.85,1.95);
\node[labelMute,anchor=east,rotate=90] at (12.40,2.80) {$I(t)$};
\node[labelMute,anchor=north] at (14.38,1.77) {time};
\draw[traceEarly] (13.16,3.38)--(14.05,2.88);
\draw[traceLate] (14.05,2.88)--(15.60,2.46);
\draw[cGray!44!black,line width=0.42pt,dash pattern=on 1.2pt off 1.2pt]
  (14.05,1.98)--(14.05,3.56);
\node[labelStd,anchor=south west] at (13.22,3.39) {2--30\,s};
\node[labelStd,anchor=north west] at (14.18,2.33) {30--300\,s};

\end{tikzpicture}%
}
\end{document}
